\section{Rings and Ideals}

\begin{exercise}\label{ex1.1}
	Consider $(1+x)\left(1-x+x^2-\cdots+(-x)^{n-1}\right)$.
\end{exercise}

\begin{exercise}\label{ex1.2}
	\begin{enumerate}[label=(\roman*),ref=\theexercise(\roman*)]\crefalias{enumi}{exercise}
		\item\label{ex1.2.1} Passing to $\left(A/\mathfrak{p}\right)[x]$. $A/\mathfrak{p}$ is an integral domain, so $\operatorname{deg}\left(\overline{f}\overline{g}\right)=\operatorname{deg}\left(\overline{f}\right)+\operatorname{deg}\left(\overline{g}\right)$. $\overline{f}\overline{g}=1$ means $a_1,\cdots,a_n\in\mathfrak{p}$. Use \Cref{1.8}.
		\item\label{ex1.2.2} $a$, $b$ are nilpotents $\Rightarrow$ $a+b$ is nilpotent. Note that $f-a_0=xf_1$ is nilpotent.
		\item\label{ex1.2.3} Let $gf=0$ and $g$ is of minimal degree. The leading coefficient $a_nb_m=0$, so $a_ng$ has less degree than $g$ and so $a_ng=0$. Same process to the second leading and following coefficients, we have $a_ig=0$ for all $i$.
		\item\label{ex1.2.4}(Gauss's lemma) Let $\mathfrak{a}$, $\mathfrak{b}$ and $\mathfrak{c}$ are generated by the coefficients of $f$, $g$ and $fg$, respectively. We need to proof $\mathfrak{c}=(1)$ $\Leftrightarrow$ $\mathfrak{ab}=(1)$. $\mathfrak{c}\subseteq\mathfrak{ab}$ by the form of coefficients in $fg$. If $\mathfrak{c}\subsetneq\mathfrak{ab}=(1)$, let $\mathfrak{m}$ be the maximal ideal containing $\mathfrak{c}$, then $\left(A/\mathfrak{m}\right)[x]$ is an integer domain. $\overline{fg}=0$ but $\overline{f}, \overline{g}\neq0$.
	\end{enumerate}
\end{exercise}

\begin{exercise}\label{ex1.3}
	Induction on the indeterminates.
\end{exercise}

\begin{exercise}\label{ex1.4}
	Pick the $y$ in \Cref{1.9} to be the indeterminate, then use \Cref{ex1.2.1}.
\end{exercise}

\begin{exercise}\label{ex1.5}
	\begin{enumerate}[label=(\roman*),ref=\theexercise(\roman*)]\crefalias{enumi}{exercise}
		\item\label{ex1.5.1} $\Leftarrow$: $fg=\sum_{i=0}^\infty\left(\sum_{j=0}^ia_jb_{i-j}\right)x^i$, so we can find $g$ such that $\sum_{j=0}^ia_jb_{i-j}=0$.
		\item\label{ex1.5.2} Similar to \Cref{ex1.2.2}. Construct a counterexample where $a_n$ have increasing nilpotencies.
		\item Pick the $y$ in \Cref{1.9} to be 1. Use \hyperref[ex1.5.1]{(i)} on $1-f$ $\Leftrightarrow$ $1-a_0$.
		\item[(iv),(v)] $A\cong A[[x]]/(x)$. Use \Cref{1.1}.
	\end{enumerate}
\end{exercise}

\begin{exercise}\label{ex1.6}
	$\mathfrak{J}$ is an ideal, so there is an idempotent $e\in\mathfrak{J}$ but $e\notin\mathfrak{R}$. Pick the $y$ in \Cref{1.9} to be 1.
\end{exercise}

\begin{exercise}\label{ex1.7}
	Note that $x\left(1-x^{n-1}\right)=0\in\mathfrak{p}$.
\end{exercise}

\begin{exercise}\label{ex1.8}
	Zorn's lemma but define the partial order reversely.
\end{exercise}

\begin{exercise}\label{ex1.9}
	\begin{enumerate}
		\item[$\Rightarrow$:] Apply \Cref{1.8} to $A/\mathfrak{a}$.
		\item[$\Leftarrow$:] Use \Cref{1.11.2}.
	\end{enumerate}
\end{exercise}

\begin{exercise}\label{ex1.10}
	Use \Cref{1.5} and \Cref{1.6.1}
\end{exercise}

\begin{exercise}\label{ex1.11}
	\begin{enumerate}[label=(\roman*),ref=\theexercise(\roman*)]\crefalias{enumi}{exercise}
		\item $x+x=(x+x)^2=x^2+2x+x^2=x+2x+x$.
		\item Proof that every element in $A/\mathfrak{p}$ is 1.
		\item\label{ex1.11.3} $x+y-xy$.
	\end{enumerate}
\end{exercise}

\begin{exercise}\label{ex1.12}
	Use \Cref{1.5}.
\end{exercise}

\begin{exercise}\label{ex1.13}
	\textcolor{magenta}{1 is a monic polynomial but it is neither irreducible nor reducible.} $A$ is an integral domain, so $\operatorname{deg}\left(fg\right)=\operatorname{deg}\left(f\right)+\operatorname{deg}\left(g\right)$. If $\mathfrak{a}=(1)$, apply the equation to $1=\sum_{f}y_ff(x_f)$.
\end{exercise}

\begin{exercise}\label{ex1.14}
	If $xy\in\mathfrak{m}$, $xy$ is zero-divisor, so are $x$ and $y$. By maximality, $\mathfrak{m}\subseteq\mathfrak{m}+(x)=\mathfrak{m}$.
\end{exercise}

\begin{exercise}\label{ex1.15}
	\begin{enumerate}
		\item[(i)] Apply \Cref{1.8} to $A/\mathfrak{a}$ to proof the second equation.
		\item[(iv)] Use $\sqrt{\mathfrak{a}\cap\mathfrak{b}}=\sqrt{\mathfrak{ab}}$.
	\end{enumerate}
\end{exercise}

\begin{exercise}\label{ex1.16}
	Irreducible polynomials relate to prime ideals.
\end{exercise}

\begin{exercise}\label{ex1.17}
	\begin{enumerate}[label=(\roman*),ref=\theexercise(\roman*)]\crefalias{enumi}{exercise}
		\item[(i)-(iv)] Use \Cref{ex1.15}.
		\addtocounter{enumi}{4}
		\item\label{ex1.17.5} Let $\left(X_{f_i}\right)_{i\in I}$ be an open cover. By \Cref{ex1.15}, $V\left(\left\{f_i\right\}_{i\in I}\right)=\varnothing$, so $\left\{f_i\right\}_{i\in I}$ generates the unit ideal. Let $1=\sum_{i\in J}g_if_i$ where $J$ is a finite subset of $I$. If $\mathfrak{p}\notin\bigcup_{i\in J}X_{f_i}$, then $f_i\in\mathfrak{p}$ for all $i\in J$.
		\item Similar to \hyperref[ex1.17.5]{(v)}, where the unit ideal shall change to the radical $\sqrt{(f)}$.
	\end{enumerate}
\end{exercise}

\begin{exercise}\label{ex1.18}
	\begin{enumerate}
		\item[(iv)] If $V\left(\mathfrak{p}_x\right)\subseteq V\left(\mathfrak{p}_y\right)$, $\mathfrak{p}_x\in\overline{V\left(\mathfrak{p}_y\right)}$. If they don't have containing relation, $\mathfrak{p}_x\in\overline{V\left(\mathfrak{p}_y\right)}$ and $\mathfrak{p}_y\in\overline{V\left(\mathfrak{p}_x\right)}$.
	\end{enumerate}
\end{exercise}

\begin{exercise}\label{ex1.19}
	Use the closed sets definition. Consider minimal prime ideals.
\end{exercise}

\begin{exercise}\label{ex1.20}
	\begin{enumerate}[label=(\roman*),ref=\theexercise(\roman*)]\crefalias{enumi}{exercise}
		\item\label{ex1.20.1} Let $\overline{Y}=Z_1\cup Z_2$, then $Z_i\cap Y=\varnothing$ for an $i\in\{1,2\}$. Thus $Z_i=\overline{Z_i\cap Y}=\varnothing$.
		\item Similar to \hyperref[ex1.20.1]{(i)}. Construct the maximal element and proof that it is irreducible.
		\item If $Y$ Hausdorff but irreducible, every pair of open sets in $Y$ intersect. Thus $Y$ is singleton.
		\item\label{ex1.20.4} Use the closed set definition.
	\end{enumerate}
\end{exercise}

\begin{exercise}\label{ex1.21}
	Simplify the sets: $\phi^{*-1}(E)=\left\{\mathfrak{q}\in Y:\phi^{-1}(\mathfrak{q})\in E\right\}$, $\phi^{*}(V(E))=\left\{\phi^{-1}(\mathfrak{q}):E\subseteq\mathfrak{q}\right\}$.
\end{exercise}

\begin{exercise}\label{ex1.22}
	Use quotient to find the form of prime ideals in $A$.
\end{exercise}

\begin{exercise}\label{ex1.23}
	\begin{enumerate}[label=(\roman*),ref=\theexercise(\roman*)]\crefalias{enumi}{exercise}
		\item $X_f=V(1-f)$.
		\item\label{ex1.23.2} $X_{f_1}\cup X_{f_2}=X_{f_1+f_2-f_1f_2}$.
		\item A closed set in compact space is compact. By \Cref{ex1.17}, it is a finite union of basic open set. Then apply \hyperref[ex1.23.2]{(ii)}.
		\item\label{ex1.23.4} By \Cref{ex1.11} every prime ideal is maximal. Pick $f\in\mathfrak{p}_x\backslash\mathfrak{p}_y$, then $X_f$ and $V(f)$ separate two points.
	\end{enumerate}
\end{exercise}

\begin{exercise}\label{ex1.24}
	$a=a\wedge1=a\wedge\left(a\vee a^\prime\right)=(a\wedge a)\vee\left(a\wedge a^\prime\right)=(a\wedge a)\vee0=a\wedge a$.
\end{exercise}

\begin{exercise}\label{ex1.25}
	Use \Cref{ex1.23.4}.
\end{exercise}

\begin{exercise}\label{ex1.26}
	$\mathfrak{m}_x\in\mu\left(U_f\right)$ $\Leftrightarrow$ $f(x)\neq0$ $\Leftrightarrow$ $f\notin\mathfrak{m}_x$
	$\Leftrightarrow$
	$\mathfrak{m}_x\in\tilde{U}_f$.
\end{exercise}

\begin{exercise}\label{ex1.27}
	If $x\neq y$, there exists a polynomial $f\in P(X)$ such that $f(x)=0$ but $f(y)\neq0$, so $\mathfrak{m}_x\neq\mathfrak{m}_y$, and so $\mu$ injective. By Nullstellensatz, every maximal ideal in $P(X)$ is of the form $\mathfrak{m}=\left(t_1-a_1,\cdots,t_n-a_n\right)$, with $\left(a_1,\cdots,a_n\right)\in X$, so $\mu$ surjective. By \Cref{ex1.26}, $\mu$ is a homeomorphism.
\end{exercise}

\begin{exercise}\label{ex1.28}
	Project to coordinate functions.
\end{exercise}